\documentclass[11pt]{amsart}

\usepackage{amsmath,amssymb,amsthm}
\usepackage{mathtools}
\usepackage{booktabs}
\usepackage{array}
\usepackage[margin=1in]{geometry}

\newtheorem{theorem}{Theorem}[section]
\newtheorem{proposition}[theorem]{Proposition}
\newtheorem{lemma}[theorem]{Lemma}
\newtheorem{corollary}[theorem]{Corollary}
\newtheorem{conjecture}[theorem]{Conjecture}

\theoremstyle{definition}
\newtheorem{definition}[theorem]{Definition}
\newtheorem{example}[theorem]{Example}
\newtheorem{remark}[theorem]{Remark}
\newtheorem{question}[theorem]{Question}

\newcommand{\N}{\mathbb{N}}
\newcommand{\Z}{\mathbb{Z}}
\newcommand{\Q}{\mathbb{Q}}
\newcommand{\fS}{\mathfrak{S}}
\newcommand{\cL}{\mathcal{L}}
\newcommand{\cD}{\mathcal{D}}
\newcommand{\cT}{\mathcal{T}}
\newcommand{\cP}{\mathcal{P}}
\newcommand{\cU}{\mathcal{U}}
\newcommand{\cN}{\mathcal{N}}
\newcommand{\NH}{\mathcal{NH}}
\newcommand{\cH}{\mathcal{H}}
\newcommand{\DAHA}{\ddot{\cH}}

\DeclareMathOperator{\id}{id}
\DeclareMathOperator{\Comp}{Comp}
\DeclareMathOperator{\Fl}{Fl}

\title{The umbral calculus with divided differences}
\author{Robin Langer}
\date{\today}

\begin{document}

\begin{abstract}
We trace the analogy between the classical umbral calculus of Rota,
its lift to the ring of symmetric functions $\Lambda$ via coalgebra
morphisms and the Hall inner product (as developed in the author's
masters thesis), and a prospective ``divided difference umbral
calculus'' on the polynomial ring $\Z[x_1,\dots,x_n]$. The column
operator construction survives: Demazure operators play the role of
generalized column operators, and key polynomials play the role of
Littlewood--Richardson type bases. However, the parametric freedom
collapses from an infinite-dimensional function space (delta series
under composition) to the Macdonald parameters~$(q,t)$, due to
the nilpotency and noncommutativity of divided differences. The
Hecke algebra, the double affine Hecke algebra, and the Cherednik
operators play the roles of the delta series, the umbral operators,
and the column operators respectively.
\end{abstract}

\maketitle

\tableofcontents

% ------------------------------------------------------------------
\section{Introduction}
% ------------------------------------------------------------------

The umbral calculus of Gian-Carlo Rota \cite{Rota,Roman} classifies
the coalgebra automorphisms of the polynomial ring $\Q[x]$ by
consideration of the algebra morphisms of the dual space
$\Q[[y]]$. Every such automorphism is determined by an invertible
formal power series $f(z)$ (a \emph{delta series}), and gives rise
to a \emph{sequence of binomial type}: a basis $\{p_n(x)\}$ for
$\Q[x]$ satisfying
\[
  p_n(x + x') = \sum_{k=0}^{n} \binom{n}{k} p_k(x)\, p_{n-k}(x').
\]

In \cite{Langer}, the author lifted this machinery from $\Q[x]$ to
the ring of symmetric functions $\Lambda$, using the coalgebra
embedding $x^n/n! \mapsto h_n(X)$ and the Hall inner product. The
result was a classification of all \emph{Littlewood--Richardson type
bases} $\{P_\lambda(X)\}$ of $\Lambda$: bases satisfying
\[
  P_\lambda(X + X') = \sum_{\mu,\nu} c^\lambda_{\mu\nu}\, P_\mu(X)\, P_\nu(X')
\]
with the same structure constants $c^\lambda_{\mu\nu}$ as the Schur
functions.

A natural question arises: what happens if we replace the
differential operator $\cD = d/dx$ not with the adjoint operators
$\partial_{h_n}$ on $\Lambda$, but with the \emph{divided difference
operators} $\partial_i$ on the polynomial ring
$\Z[x_1,\dots,x_n]$? This note traces the analogy, identifies what
survives, what breaks, and what replaces the broken parts.

% ------------------------------------------------------------------
\section{Three levels of the umbral calculus}
% ------------------------------------------------------------------

% ------------------------------------------------------------------
\subsection{Level 1: The classical setting}\label{sec:classical}
% ------------------------------------------------------------------

The polynomial ring $V = \Q[x]$ has a distinguished Hamel basis
$\{1, x, x^2/2!, x^3/3!, \dots\}$ and a coalgebra structure given
by:
\[
  \Delta\, p(x) = p(x + x').
\]
The dual space $V^* = \Q[[y]]$ has basis $\{1, y, y^2, \dots\}$,
and the pairing between $V^*$ and $V$ is:
\[
  (y^k, p(x)) = \left[\frac{x^k}{k!}\right] p(x) = \cL \cD^k\, p(x)
\]
where $\cD = d/dx$ is the differential operator and $\cL$ is the
constant term operator.

A \emph{delta series} is a formal power series $f(y) \in \Q[[y]]$
with $f(0) = 0$ and $f'(0) \neq 0$, i.e., one which has a
compositional inverse $g(y)$ satisfying $f(g(y)) = g(f(y)) = y$.

\begin{proposition}[{\cite[Prop.~1.2]{Langer}}]
  For every delta series $f(y)$ with compositional inverse $g(y)$,
  the \emph{umbral operator}
  \[
    U_f = \sum_n \frac{x^n}{n!}\, \cL\, f^n(\cD)
  \]
  is the unique operator on $\Q[x]$ satisfying:
  \begin{enumerate}
    \item $\cD \circ U_f = U_f \circ g(\cD)$,
    \item $\cL \circ U_f = \cL$.
  \end{enumerate}
\end{proposition}

The group of delta series under composition --- the
\emph{Jabotinsky group} (or Riordan group) --- is
infinite-dimensional. This is the source of the parametric freedom
in the classical umbral calculus: there are infinitely many
sequences of binomial type, one for each delta series.

% ------------------------------------------------------------------
\subsection{Level 2: Symmetric functions}\label{sec:symmetric}
% ------------------------------------------------------------------

The coalgebra embedding
\[
  \frac{x^n}{n!} \mapsto h_n(X)
\]
lifts the coproduct from $\Q[x]$ to $\Lambda$, where
$\Delta\, f(X) = f(X + X')$ is given by plethystic addition of
alphabets. The Hall inner product
\[
  \langle m_\lambda(Y), h_\mu(X) \rangle = \delta_{\lambda\mu}
\]
replaces the pairing $(y^k, p(x)) = \cL \cD^k p(x)$. The role of
the differential operator $\cD$ is played by the adjoint operators
$\partial_{h_n}$, defined by:
\[
  \langle \partial_f[g(Y)], h(X) \rangle
  = \langle g(Y),\, f(X)\, h(X) \rangle.
\]

The \emph{translation operator}
\[
  \cT(z)[f(X)] = f(X + z) = \sum_n z^n\, \partial_{h_n}
\]
and the \emph{multiplication operator}
\[
  \cP(z)[f(X)] = \Omega(Xz)\, f(X) = \sum_n z^n\, h_n(X) \cdot f(X)
\]
combine to form the \emph{column operator}:
\[
  C(z) = \cP(-z)\, \cT(1/z).
\]

\begin{proposition}[{\cite[Prop.~1.4]{Langer}}]
  The Schur function has the column operator expression:
  \[
    s_\lambda(X) = C_{\lambda'_1} \circ C_{\lambda'_2} \circ \cdots
    \circ C_{\lambda'_m}[1].
  \]
\end{proposition}

The umbral operator on $\Lambda$ is defined by
$U_f[\Omega(Xz)] = \Omega_f(Xz) = \prod_x 1/(1 - f(z)x)$, extended
multiplicatively. The generalized column operator is:
\[
  C_f(z) = U_f \circ C(z) \circ U_g
\]
and the resulting bases
\[
  P_\lambda(X) = C^f_{\lambda'_1} \circ C^f_{\lambda'_2} \circ
  \cdots \circ C^f_{\lambda'_m}[1] = U_f[s_\lambda(X)]
  = \det(r_{\lambda_i + j - i}(X))
\]
are of Littlewood--Richardson type, with dual bases
$Q_\lambda(Y) = U_g^*[s_\lambda(Y)]$ satisfying:
\[
  \sum_\lambda P_\lambda(X)\, Q_\lambda(Y) = \Omega(XY)
  = \prod_{\substack{x \in X \\ y \in Y}} \frac{1}{1 - xy}.
\]

The parametric freedom is the same as in the classical case: every
delta series $f(y)$ gives a Littlewood--Richardson type basis.

% ------------------------------------------------------------------
\subsection{Level 3: Divided differences}\label{sec:divided}
% ------------------------------------------------------------------

Now consider the polynomial ring $\Z[x_1, \dots, x_n]$ with the
\emph{divided difference operators}:
\[
  \partial_i(f) = \frac{f - s_i(f)}{x_i - x_{i+1}},
  \qquad i = 1, \dots, n-1
\]
where $s_i$ is the simple transposition swapping $x_i$ and
$x_{i+1}$. These generate the \emph{nil-Coxeter algebra} $\cN_n$
with relations:
\begin{align}
  \partial_i^2 &= 0, \label{eq:nilpotent} \\
  \partial_i \partial_j &= \partial_j \partial_i
  \quad \text{for } |i - j| \geq 2, \\
  \partial_i \partial_{i+1} \partial_i &= \partial_{i+1} \partial_i
  \partial_{i+1} \quad \text{(braid relation)}.
\end{align}

The \emph{Demazure operators} (isobaric divided differences) are:
\[
  \pi_i(f) = \partial_i(x_i\, f)
  = \frac{x_i\, f - x_{i+1}\, s_i(f)}{x_i - x_{i+1}}.
\]
These satisfy $\pi_i^2 = \pi_i$ (idempotence) and the same braid
relations as the $\partial_i$.

% ------------------------------------------------------------------
\section{The dictionary}
% ------------------------------------------------------------------

The following table records the analogy across the three levels.

\medskip
\begin{center}
\renewcommand{\arraystretch}{1.4}
\begin{tabular}{>{\raggedright}p{3.2cm} >{\raggedright}p{4cm}
  >{\raggedright\arraybackslash}p{4.8cm}}
\toprule
\textbf{Classical} $\Q[x]$ &
\textbf{Symmetric} $\Lambda$ &
\textbf{Divided difference} $\Z[x_1,\dots,x_n]$ \\
\midrule
$\cD = d/dx$ &
$\partial_{h_n}$ &
$\partial_i = (f - s_i f)/(x_i - x_{i+1})$ \\[4pt]
$\cD^n$ &
$\partial_{h_\lambda}$ &
$\partial_w = \partial_{i_1} \cdots \partial_{i_k}$ (BGG) \\[4pt]
$x^n / n!$ &
$h_n(X)$ &
$\mathfrak{S}_w(x)$ (Schubert polynomials) \\[4pt]
$\cL \cD^k$ extracts $[x^k/k!]$ &
$\langle m_\lambda(Y), -\rangle$ &
$\partial_{w_0}$ as top extraction \\[4pt]
$\Delta p = p(x + x')$ &
$\Delta f = f(X + X')$ &
no clean analogue (\S\ref{sec:coproduct}) \\[4pt]
$\cT = \sum x^n/n!\; \cL \cD^n$ (Taylor) &
$\cT(z) = \sum z^n \partial_{h_n}$ &
$f = \sum_w \partial_w(f)|_0 \cdot \mathfrak{S}_w$ \\[4pt]
$\cP(z) = e^{xz} \times$ &
$\cP(z) = \Omega(Xz) \times$ &
--- \\[4pt]
$C(z) = \cP(-z)\cT(1/z)$ &
same &
$\pi_i = x_i \partial_i$ (Demazure) \\[4pt]
$p_n = U_f[x^n/n!]$ &
$P_\lambda = U_f[s_\lambda]$ &
$\kappa_\alpha = \pi_{i_1} \cdots \pi_{i_k}[x^\lambda]$ \\[4pt]
$\Q[\cD]$ (commutative) &
same &
$\cN_n$ (nil-Coxeter, dim $= n!$) \\[4pt]
Delta series $f(y)$ &
same &
Hecke parameter $q$ \\[4pt]
Jabotinsky group ($\infty$-dim) &
same &
$\mathrm{Aut}(\cN_n)$ (finite) \\[4pt]
$\Omega(XY) = \sum s_\lambda s_\lambda$ &
same &
$\prod_{i+j \leq n+1}(x_i + y_j) = \sum \mathfrak{S}_w \mathfrak{S}_{w_0 w}$ \\
\bottomrule
\end{tabular}
\end{center}

% ------------------------------------------------------------------
\section{What survives: the column operator construction}
% ------------------------------------------------------------------

The column operator construction is the part of the umbral calculus
that transfers most cleanly to divided differences.

\subsection{Column operators and Demazure operators}

In \S\ref{sec:symmetric}, the column operator $C(z) = \cP(-z)\cT(1/z)$
combines multiplication by a generating function with translation.
The Demazure operator $\pi_i = x_i \partial_i$ has exactly this
structure: it is multiplication by $x_i$ followed by the divided
difference $\partial_i$.

The Schur function is built by iterated column operators:
\[
  s_\lambda(X) = C_{\lambda'_1} \circ C_{\lambda'_2} \circ \cdots
  \circ C_{\lambda'_m}[1].
\]
The key polynomial is built by iterated Demazure operators:
\[
  \kappa_\alpha = \pi_{i_1} \circ \pi_{i_2} \circ \cdots \circ
  \pi_{i_k}[x^\lambda]
\]
where $s_{i_1} \cdots s_{i_k}$ is a reduced decomposition and
$\lambda$ is the partition obtained by sorting $\alpha$.

\subsection{Generalized column operators}

In the symmetric function setting, the generalized column operator
$C_f(z) = U_f \circ C(z) \circ U_g$ is conjugation of the Schur
column operator by the umbral operators. The resulting bases
$P_\lambda = C^f_{\lambda'_1} \circ \cdots \circ C^f_{\lambda'_m}[1]$
are of Littlewood--Richardson type.

The divided difference analogue of this conjugation replaces $\pi_i$
with a \emph{deformed Demazure operator}. The natural deformation
is provided by the \textbf{Hecke algebra} $\cH_n(t)$, which
replaces $s_i^2 = 1$ with:
\begin{equation}\label{eq:hecke}
  T_i^2 = (t - 1)\, T_i + t.
\end{equation}
The Hecke generator acts on the polynomial ring via the
\emph{Demazure--Lusztig operator}:
\begin{equation}\label{eq:DL}
  T_i = (t - 1)\, \pi_i + s_i.
\end{equation}
One may verify \eqref{eq:hecke} directly using $\pi_i^2 = \pi_i$,
$s_i^2 = \id$, $s_i \pi_i = \pi_i$, and
$\pi_i s_i = \id + s_i - \pi_i$:
\begin{align*}
  T_i^2 &= (t-1)^2 \pi_i + (t-1)(\pi_i s_i + s_i \pi_i) + \id \\
  &= (t-1)^2 \pi_i + (t-1)(\id + s_i) + \id \\
  &= (t-1)^2 \pi_i + (t-1) s_i + t \\
  &= (t-1) T_i + t. \qquad \checkmark
\end{align*}

The formula \eqref{eq:DL} makes the role of the parameter $t$
transparent: it controls the mixture of the Demazure projection
$\pi_i$ and the reflection $s_i$.

\medskip
\begin{center}
\renewcommand{\arraystretch}{1.3}
\begin{tabular}{lll}
\toprule
\textbf{Parameter} & \textbf{Operator} & \textbf{Polynomials} \\
\midrule
$t = 1$ & $T_i = s_i$ (reflection) &
  Schur / symmetric functions \\
$t = 0$ & $T_i = s_i - \pi_i$ (0-Hecke) &
  key polynomials / atoms \\
$t = -1$ & $T_i = -2\pi_i + s_i$ &
  --- \\
general $t$ & $T_i = (t-1)\pi_i + s_i$ &
  Hall--Littlewood / Macdonald \\
\bottomrule
\end{tabular}
\end{center}

\begin{remark}\label{rmk:0hecke}
  At $t = 0$, the relation becomes $T_i^2 = -T_i$, so $T_i + 1$
  is idempotent: $(T_i + 1)^2 = T_i + 1$. This is the
  \textbf{0-Hecke algebra}, which is the algebra of Demazure
  operators. Note that the \emph{nil-Coxeter algebra}
  ($\partial_i^2 = 0$) is a \emph{different} quotient of the braid
  group algebra --- it is \textbf{not} a specialization of the Hecke
  algebra at any value of $t$. The divided difference $\partial_i$
  requires both $s_i$ and the polynomial ring for its definition:
  $\partial_i = (1 - s_i)/(x_i - x_{i+1})$. It is an element of
  the \emph{degenerate affine Hecke algebra}, not the finite Hecke
  algebra.
\end{remark}

\subsection{Anti-Demazure operators and atoms}

The anti-Demazure operator
\[
  \bar{\pi}_i = \pi_i - \id
\]
satisfies $\pi_i + \bar{\pi}_i = \id$, so the two operators are
complementary projections. Since $\pi_i = \id + \bar{\pi}_i$,
expanding a product of Demazure operators yields:
\[
  \kappa_\alpha = \sum_{\beta \leq \alpha} A_\beta
\]
where the \emph{Demazure atoms}
$A_\alpha = \bar{\pi}_{i_1} \circ \cdots \circ
\bar{\pi}_{i_k}[x^\lambda]$ are obtained by using
$\bar{\pi}$ at every step. This decomposition is the direct
analogue of the expansion $P_\lambda(X) = \sum_{\nu \subseteq
\lambda} \square_{\lambda\nu}\, s_\nu(X)$ in
\S\ref{sec:symmetric}, where the Schur expansion coefficients
arise from the ``identity vs.\ umbral'' choice at each step.

% ------------------------------------------------------------------
\section{What breaks: three structural obstructions}
% ------------------------------------------------------------------

% ------------------------------------------------------------------
\subsection{The coproduct does not lift}\label{sec:coproduct}
% ------------------------------------------------------------------

The coproduct $\Delta f(X) = f(X + X')$ on $\Lambda$, given by
plethystic addition of alphabets, is the engine of the classical
theory. It gives rise to the binomial-type identity
$p_n(x + x') = \sum \binom{n}{k} p_k\, p_{n-k}$ and its symmetric
function generalization
$P_\lambda(X + X') = \sum c^\lambda_{\mu\nu} P_\mu P_\nu$.

The polynomial ring $\Z[x_1, \dots, x_n]$ has variable-wise
addition $\Delta p(x) = p(x_1 + y_1, \dots, x_n + y_n)$, but
Schubert polynomials do \emph{not} satisfy a binomial-type property
under this coproduct. That is, there is no identity of the form:
\[
  \mathfrak{S}_w(x + y) = \sum_{u,v} c^w_{uv}\, \mathfrak{S}_u(x)\,
  \mathfrak{S}_v(y).
\]
The ``Cauchy identity'' for Schubert polynomials is instead a
\emph{finite product}:
\[
  \prod_{i + j \leq n+1} (x_i + y_j) = \sum_{w \in \fS_n}
  \mathfrak{S}_w(x)\, \mathfrak{S}_{w_0 w}(y)
\]
which represents the class of a point in the cohomology ring
$H^*(\mathrm{Fl}(n) \times \mathrm{Fl}(n))$. This is Poincar\'e
duality, not a Cauchy kernel --- it is a polynomial, not a formal
power series.

\begin{remark}
  The theory of \emph{back-stable Schubert polynomials}
  (Lam--Lee--Shimozono \cite{LLS}) embeds Schubert polynomials into
  a larger ring that \emph{does} have a good coproduct, partially
  recovering the coalgebra structure. This may be the right setting
  for a full divided-difference umbral calculus.
\end{remark}

% ------------------------------------------------------------------
\subsection{The operator algebra is noncommutative}
% ------------------------------------------------------------------

The umbral operator $U_f = \sum x^n/n!\; \cL f^n(\cD)$ relies on
$\cD$ generating a commutative algebra $\Q[\cD]$. The delta series
$f(y)$ and its compositional inverse $g(y)$ live in the group of
invertible formal power series under composition --- the Jabotinsky
(or Riordan) group, which is infinite-dimensional.

The divided differences $\partial_1, \dots, \partial_{n-1}$ generate
the nil-Coxeter algebra $\cN_n$, which is:
\begin{itemize}
  \item \textbf{noncommutative} (the braid relation is not
    commutativity),
  \item \textbf{finite-dimensional} ($\dim \cN_n = n!$, with basis
    $\{\partial_w : w \in \fS_n\}$),
  \item \textbf{nilpotent} ($\partial_i^2 = 0$).
\end{itemize}

The finiteness and noncommutativity mean that ``formal power series
in $\partial$'' do not form a commutative ring, and compositional
inversion of delta series has no direct analogue. The entire
Jabotinsky group collapses to the automorphism group of $\cN_n$,
which is finite.

% ------------------------------------------------------------------
\subsection{Lowering happens in multiple directions}
% ------------------------------------------------------------------

In $\Q[x]$, the operator $\cD$ lowers degree by~1 along a single
axis. Taylor's formula
$\mathcal{I} = \sum x^n/n!\; \cL \cD^n$ expresses the identity as
a sum over $\N$.

With divided differences, there are $n - 1$ operators
$\partial_1, \dots, \partial_{n-1}$, each lowering degree by~1
but in a different root direction. The Taylor formula becomes the
\emph{Schubert expansion}:
\[
  f = \sum_{w \in \fS_n} \partial_w(f)\big|_{x = 0} \cdot
  \mathfrak{S}_w(x)
\]
which is a sum over the symmetric group $\fS_n$ rather than over
$\N$. The multi-directionality is intrinsic to the theory and
cannot be reduced to a single operator.

% ------------------------------------------------------------------
\section{What replaces the broken parts}
% ------------------------------------------------------------------

% ------------------------------------------------------------------
\subsection{The Hecke parameter replaces delta series}
% ------------------------------------------------------------------

In the classical theory, the parametric family of umbral operators
is indexed by delta series $f(y)$ --- an infinite-dimensional space.
In the divided difference setting, the natural deformation is the
Hecke algebra $\cH_n(t)$, which replaces the symmetric group
relation $s_i^2 = 1$ with \eqref{eq:hecke}.

As shown in Remark~\ref{rmk:0hecke}, $t = 0$ gives the 0-Hecke
algebra (the algebra of Demazure operators, with key polynomials
and atoms), while $t = 1$ recovers the symmetric group (with Schur
polynomials). At generic $t$, the Demazure--Lusztig operators
generate Hall--Littlewood polynomials.

To obtain the full Macdonald polynomials, one needs the second
parameter $q$, which arises from the \emph{double affine Hecke
algebra} (see \S\ref{sec:DAHA} below).

\begin{remark}
  The collapse from an infinite-dimensional parameter space (delta
  series) to a two-dimensional one (Macdonald's $(q,t)$) reflects
  the collapse from $\Q[\cD] \cong \Q[[y]]$
  (infinite-dimensional) to the finite-dimensional operator
  algebras generated by divided differences. The braid group
  algebra admits only finitely many interesting quotients, so only
  finitely many deformation parameters survive.
\end{remark}

% ------------------------------------------------------------------
\subsection{The DAHA and Cherednik operators}\label{sec:DAHA}
% ------------------------------------------------------------------

The single Hecke parameter $t$ gives Hall--Littlewood polynomials.
To obtain the full Macdonald polynomials $P_\lambda(x; q, t)$, one
needs the \emph{double affine Hecke algebra} (DAHA), introduced by
Cherednik \cite{Cherednik}. The DAHA $\DAHA_n(q,t)$ is generated by
three families of operators:
\begin{itemize}
  \item $T_1, \dots, T_{n-1}$ (Hecke generators, parameter $t$),
  \item $x_1, \dots, x_n$ (polynomial multiplication),
  \item $Y_1, \dots, Y_n$ (Cherednik operators),
\end{itemize}
subject to the cross-relations:
\begin{align}
  T_i\, x_i\, T_i &= t\, x_{i+1}, \label{eq:daha-x} \\
  T_i\, Y_i\, T_i &= Y_{i+1}. \label{eq:daha-y}
\end{align}
The Cherednik operators commute:
\begin{equation}\label{eq:cherednik-commute}
  Y_i\, Y_j = Y_j\, Y_i \qquad \text{for all } i, j.
\end{equation}

The \emph{nonsymmetric Macdonald polynomials}
$E_\alpha(x; q, t)$ are the simultaneous eigenfunctions of the
Cherednik operators:
\[
  Y_i\, E_\alpha = \xi_i(\alpha; q, t) \cdot E_\alpha
\]
where $\xi_i(\alpha; q, t)$ are explicit eigenvalues depending on
the composition $\alpha$ and the parameters.

\subsubsection{The Cherednik operators as umbral operators}

In \cite{Langer}, the umbral operator $U_f$ is characterized by the
intertwining property:
\[
  \cD \circ U_f = U_f \circ g(\cD).
\]
This says that $U_f$ \emph{conjugates} the derivative into a new
operator $g(\cD)$. The Cherednik operators satisfy an analogous
intertwining relation \eqref{eq:daha-y}: the Hecke generator $T_i$
permutes the Cherednik operators, just as $U_f$ transforms the
derivative.

More precisely, the analogy is:

\medskip
\begin{center}
\renewcommand{\arraystretch}{1.3}
\begin{tabular}{>{\raggedright}p{5.5cm}
  >{\raggedright\arraybackslash}p{5.5cm}}
\toprule
\textbf{Umbral calculus on $\Lambda$} &
\textbf{DAHA} \\
\midrule
Umbral operator $U_f$ &
Cherednik operator $Y_i$ \\[3pt]
$\cD \circ U_f = U_f \circ g(\cD)$ &
$T_i\, Y_i\, T_i = Y_{i+1}$ \\[3pt]
Delta series $f(y)$ &
Parameters $(q, t)$ \\[3pt]
Jabotinsky group ($\infty$-dim) &
Two parameters \\[3pt]
$U_f[s_\lambda] = P_\lambda$ &
$Y_i\, E_\alpha = \xi_i \cdot E_\alpha$ \\[3pt]
$\Omega(XY) = \sum P_\lambda Q_\lambda$ &
Macdonald Cauchy identity \\
\bottomrule
\end{tabular}
\end{center}

\subsubsection{Intertwining operators as column operators}

The closest analogue of the column operators
$C_f(z) = U_f \circ C(z) \circ U_g$ from \cite{Langer} is the
family of \emph{Knop--Sahi intertwining operators} $\phi_i$ in the
DAHA \cite{Knop,Sahi}. These satisfy:
\begin{equation}\label{eq:intertwine}
  Y_i\, \phi_i = \phi_i\, Y_{i+1}
\end{equation}
and build up the Macdonald polynomials step by step:
\[
  E_{s_i \cdot \alpha} \propto \phi_i(E_\alpha)
\]
starting from a ``base'' polynomial. This is directly analogous to
the column operator construction
\[
  s_\lambda(X) = C_{\lambda'_1} \circ C_{\lambda'_2} \circ \cdots
  \circ C_{\lambda'_m}[1]
\]
from \cite[Prop.~1.4]{Langer}: the intertwining operators are the
``generalized column operators'' of the Hecke world.

\subsubsection{The Macdonald involution as compositional inverse}

In \cite{Langer}, the umbral operator $U_f$ and its adjoint
$U_g^*$ are related by the compositional inverse $g = f^{-1}$:
the LR-type basis $P_\lambda = U_f[s_\lambda]$ has dual
$Q_\lambda = U_g^*[s_\lambda]$, and the duality is controlled by
$f \leftrightarrow g$.

In the Macdonald setting, the analogous duality is the
\emph{Macdonald involution}: the symmetry that sends
$P_\lambda(x; q, t)$ to $Q_\lambda(x; q, t)$. This involution
involves interchanging the roles of the two parameters (roughly
$q \leftrightarrow t^{-1}$, up to normalizations). The
compositional inverse $g = f^{-1}$ of the delta series --- which
drives the entire duality theory in \cite{Langer} --- becomes the
$(q, t)$-duality of Macdonald polynomials.

% ------------------------------------------------------------------
\subsection{The nilHecke algebra replaces the operator algebra}
% ------------------------------------------------------------------

The \emph{nilHecke algebra} $\NH_n$ is the smash product
$\Z[x_1, \dots, x_n] \rtimes \cN_n$, with the cross-relation:
\begin{equation}\label{eq:nilhecke}
  \partial_i\, x_i = 1 + x_{i+1}\, \partial_i.
\end{equation}
This is the divided-difference Leibniz rule, and it plays the role
of the commutation relation
$\cD \circ U_f = U_f \circ g(\cD)$ from
\cite[Prop.~1.2]{Langer}. The nilHecke algebra is the single
algebra in which both ``multiplication'' ($\cP$) and
``translation'' ($\cT$) operators naturally live.

\begin{remark}
  In the symmetric function setting, the translation and
  multiplication operators satisfy the commutation relation
  (Lemma~1.1 of \cite{Langer}):
  \[
    \cT(v)\, \cP(w) = \Omega(vw)\, \cP(w)\, \cT(v).
  \]
  The nilHecke relation \eqref{eq:nilhecke} is the pointwise,
  operator-level version of this: $\partial_i$ (a single step of
  translation) and $x_i$ (multiplication by a single variable)
  do not commute, and their commutator is controlled by $s_i$.
\end{remark}

% ------------------------------------------------------------------
\subsection{The Yang--Baxter equation replaces the binomial theorem}
% ------------------------------------------------------------------

The binomial-type identity
$p_n(x + x') = \sum \binom{n}{k} p_k(x)\, p_{n-k}(x')$ is the
defining property of sequences of binomial type. In the divided
difference setting, the structural identity that replaces it is the
\emph{Yang--Baxter equation}.

Following Fomin and Kirillov \cite{FK}, define the
\emph{Yang--Baxter exponential} in the nil-Coxeter algebra:
\[
  \mathfrak{E}(z) = \prod_{i} (1 + z\, u_i)
\]
where $u_i$ are the nil-Coxeter generators. This satisfies:
\[
  (1 + z\, u_i)(1 + z\, u_{i+1})(1 + z\, u_i)
  = (1 + z\, u_{i+1})(1 + z\, u_i)(1 + z\, u_{i+1}).
\]

The Yang--Baxter equation is the ``binomial theorem'' for divided
differences: it governs how the exponential factors when the
``alphabet'' is split. Just as the binomial theorem for $p_n$ is
equivalent to $\Delta p_n = \sum \binom{n}{k} p_k \otimes p_{n-k}$
being a coalgebra map, the Yang--Baxter equation is equivalent to a
\emph{quasi-triangularity} condition on the nil-Coxeter algebra.

% ------------------------------------------------------------------
\subsection{Poincar\'e duality replaces the Cauchy kernel}
% ------------------------------------------------------------------

The Cauchy kernel
$\Omega(XY) = \prod 1/(1-xy) = \sum_\lambda P_\lambda(X)\, Q_\lambda(Y)$
characterizes dual bases in $\Lambda$ with respect to the Hall inner
product. In the divided difference setting, the ``Cauchy identity''
is the finite product:
\[
  \prod_{i + j \leq n+1} (x_i + y_j) = \sum_{w \in \fS_n}
  \mathfrak{S}_w(x)\, \mathfrak{S}_{w_0 w}(y).
\]
The dual of $\mathfrak{S}_w$ is $\mathfrak{S}_{w_0 w}$, where
$\ell(w) + \ell(w_0 w) = \ell(w_0) = \binom{n}{2}$, reflecting
Poincar\'e duality on the flag variety $\mathrm{Fl}(n)$.

\medskip

\begin{center}
\renewcommand{\arraystretch}{1.3}
\begin{tabular}{>{\raggedright}p{5cm}
  >{\raggedright\arraybackslash}p{6cm}}
\toprule
\textbf{Symmetric functions} & \textbf{Schubert calculus} \\
\midrule
$\mathrm{Rep}(GL_n)$ & $H^*(\mathrm{Fl}(n))$ \\
Cauchy kernel (power series) & finite product (polynomial) \\
$s_\lambda$ self-dual & $\mathfrak{S}_{w_0 w}$ dual to
  $\mathfrak{S}_w$ \\
$\sum_\lambda s_\lambda s_\lambda = \Omega$ &
  $\sum_w \mathfrak{S}_w \mathfrak{S}_{w_0 w} = \prod(x_i + y_j)$ \\
LR coefficients & Schubert structure constants \\
\bottomrule
\end{tabular}
\end{center}

% ------------------------------------------------------------------
\section{Key polynomials and Demazure atoms as an LR-type pair}
% ------------------------------------------------------------------

The Cauchy identity for key polynomials and Demazure atoms is:
\[
  \prod_{i,j=1}^{n} \frac{1}{1 - x_i y_j} = \sum_{\alpha \in
  \Comp_n} \kappa_\alpha(x_1, \dots, x_n) \cdot
  A_{w_0 \cdot \alpha}(y_1, \dots, y_n)
\]
where $w_0 \cdot \alpha$ reverses the composition. This is the
closest analogue of the symmetric function identity
$\sum_\lambda P_\lambda(X)\, Q_\lambda(Y) = \Omega(XY)$ from
\cite{Langer}: the key polynomials play the role of $P_\lambda$
and the Demazure atoms play the role of $Q_\lambda$.

The decomposition
$\kappa_\alpha = \sum_{\beta \leq \alpha} A_\beta$
is upper-triangular, making key polynomials and atoms dual bases
--- exactly as the Littlewood--Richardson type bases $P_\lambda$
and their duals $Q_\lambda$ are related by an upper-triangular
change of basis from the Schur functions.

% ------------------------------------------------------------------
\section{The collapse of parametric freedom}
% ------------------------------------------------------------------

\begin{question}
  In the classical umbral calculus, the Jabotinsky group of delta
  series parametrizes all sequences of binomial type. What is
  the ``Jabotinsky group'' for divided differences?
\end{question}

The answer appears to be: \emph{essentially trivial}. The
automorphism group of the nil-Coxeter algebra $\cN_n$ is generated
by the diagram automorphisms of the Dynkin diagram $A_{n-1}$ (i.e.,
the map $\partial_i \mapsto \partial_{n-i}$, corresponding to
$w \mapsto w_0 w w_0$). This is a finite group, in stark contrast
to the infinite-dimensional Jabotinsky group.

The Hecke deformation $s_i^2 = 1 \;\rightsquigarrow\; T_i^2
= (t-1)T_i + t$ provides a one-parameter family, and the
DAHA (\S\ref{sec:DAHA}) extends this to the two-parameter
Macdonald $(q,t)$-family. But these are \emph{finitely many}
parameters, not the infinite-dimensional family of delta series.

\begin{conjecture}
  The Macdonald $(q,t)$-parameters exhaust the deformation freedom
  available when the derivative is replaced by divided differences
  in the umbral calculus. More precisely, any ``divided-difference
  umbral operator'' that produces a basis of $\Z[x_1, \dots, x_n]$
  with Littlewood--Richardson type structure constants is
  equivalent, up to automorphisms of $\cN_n$, to a specialization
  of the Macdonald construction.
\end{conjecture}

% ------------------------------------------------------------------
\section{Directions}
% ------------------------------------------------------------------

\begin{enumerate}
  \item \textbf{Back-stable Schubert polynomials.} The work of
    Lam--Lee--Shimozono \cite{LLS} embeds Schubert polynomials into
    a ring with a good coproduct. Does this restore enough coalgebra
    structure to support a full umbral calculus?

  \item \textbf{Noncommutative symmetric functions.} The theory of
    Gelfand--Krob--Lascoux--Leclerc--Thibon--Retakh \cite{GKLLTR}
    provides a Hopf algebra of noncommutative symmetric functions
    with divided differences built in. Is this the ``right'' home
    for the divided-difference umbral calculus?

  \item \textbf{The nilHecke algebra as a bialgebra.} The nilHecke
    algebra $\NH_n$ has both a product (algebra structure) and a
    coproduct (induced from the inclusion
    $\fS_m \times \fS_k \hookrightarrow \fS_{m+k}$). Does the
    resulting bialgebra structure support a notion of
    ``binomial type''?

  \item \textbf{Elliptic deformations.} Beyond the Hecke parameter
    $q$ and the Macdonald parameters $(q,t)$, elliptic
    deformations (Ruijsenaars, Rains) provide further parameters.
    Are these additional ``delta series'' in the divided-difference
    setting?

  \item \textbf{Categorification.} The nilHecke algebra
    categorifies the quantum group $U_q(\mathfrak{sl}_n)$ via
    Khovanov--Lauda--Rouquier algebras. Does the umbral calculus
    have a categorical shadow?
\end{enumerate}

% ------------------------------------------------------------------
% References
% ------------------------------------------------------------------

\begin{thebibliography}{99}

\bibitem{Cherednik}
I.~Cherednik,
\emph{Double affine Hecke algebras},
London Math.\ Soc.\ Lecture Note Series \textbf{319},
Cambridge University Press, 2005.

\bibitem{FK}
S.~Fomin and A.~N.~Kirillov,
\emph{The Yang--Baxter equation, symmetric functions, and Schubert
  polynomials},
Discrete Math. \textbf{153} (1996), 123--143.

\bibitem{GKLLTR}
I.~M.~Gelfand, D.~Krob, A.~Lascoux, B.~Leclerc, V.~Retakh, and
J.-Y.~Thibon,
\emph{Noncommutative symmetric functions},
Adv.\ Math. \textbf{112} (1995), 218--348.

\bibitem{Knop}
F.~Knop,
\emph{Integrality of two variable Kostka functions},
J.~Reine Angew.\ Math. \textbf{482} (1997), 177--189.

\bibitem{Langer}
R.~Langer,
\emph{Symmetric functions and Macdonald polynomials},
Masters thesis, University of Melbourne, 2008.

\bibitem{LLS}
T.~Lam, S.~J.~Lee, and M.~Shimozono,
\emph{Back stable Schubert calculus},
Compositio Math. \textbf{157} (2021), 883--962.

\bibitem{Macdonald}
I.~G.~Macdonald,
\emph{Symmetric functions and Hall polynomials},
2nd ed., Oxford University Press, 1995.

\bibitem{Reiner-Shimozono}
V.~Reiner and M.~Shimozono,
\emph{Key polynomials and a flagged Littlewood--Richardson rule},
J.~Combin.\ Theory Ser.~A \textbf{70} (1995), 107--143.

\bibitem{Roman}
S.~Roman,
\emph{The Umbral Calculus},
Academic Press, 1984.

\bibitem{Rota}
G.-C.~Rota,
\emph{Finite operator calculus},
Academic Press, 1975.

\bibitem{Sahi}
S.~Sahi,
\emph{Interpolation, integrality, and a generalization of
  Macdonald's polynomials},
Internat.\ Math.\ Res.\ Notices \textbf{1996}, no.~10, 457--471.

\end{thebibliography}

\end{document}
